In key exchange protocols, the derived key should be bound to the session. In TLS 1.3, this binding happens rather late, which has been identified as a significant source of complexity for analysing the TLS 1.3 key schedule (ASIACRYPT 2022).

We argue that the extension of TLS 1.3 would be a natural point to adapt early binding. While the computational costs of an early binding to both the ciphertext and public key of the KEM as well as the key shares of the DH would be negligible, we show the significant difference in complexity of the analysis by providing two formally verified proofs of a key exchange, once with early and once with late binding. We model the key exchange in the state separating proofs framework (SSP, ASIACRYPT 2018), and use the recently published Domino tool to verify our proofs.
